\hypertarget{target}{%
\subsubsection{Target}\label{target}}

Human RBX1 (UniProt P62877, 108 aa). The exact sequence provided to
competitors is in \texttt{data/target/rbx1.fasta}. The recombinant
construct used for expression at the Adaptyv Foundry is the full-length
P62877 with a C-terminal purification tag (vendor catalogue pending).

\hypertarget{submission-and-selection}{%
\subsubsection{Submission and
selection}\label{submission-and-selection}}

198 teams submitted 12,707 designs through the GEM × Adaptyv 2026
portal. Each team was capped at 100 ranked sequences (≤250 aa, ≤75\%
identity to known proteins, standard amino acids only). An expert panel
from the GEM workshop selected 322 designs across 48 distinct teams for
wet-lab validation.

\hypertarget{wet-lab-pipeline}{%
\subsubsection{Wet-lab pipeline}\label{wet-lab-pipeline}}

Designs were synthesised, expressed, and screened at the Adaptyv
Foundry. Expression was assessed on Coomassie-stained gels and BLI
loading checks. Binding was measured by surface plasmon resonance
(Carterra LSA) with 2--3 replicates per design (943 total SPR records
across 318 designs with replicates; three designs additionally have BLI
replicates). Per-replicate KD fits were averaged after curated-
replicate filtering at ProteinBase.

\hypertarget{binding-strength-bins}{%
\subsubsection{Binding-strength bins}\label{binding-strength-bins}}

ProteinBase classifies each design into the standard hierarchy:
\texttt{Strong\ \textgreater{}\ Medium\ \textgreater{}\ Weak\ \textgreater{}\ Potential\ binder\ \textgreater{}\ Non-binder\ \textgreater{}\ No\ expression\ \textgreater{}\ Unknown}.
We treat \texttt{binding\_strength\ ∈\ \{Strong,\ Medium,\ Weak\}} as a
binder for headline counts (\texttt{is\_binder} in
\texttt{data/designs.csv}). The \emph{Strong} bin maps to KD \textless{}
100 nM.

\hypertarget{in-silico}{%
\subsubsection{In silico}\label{in-silico}}

Every design was re-folded by ProteinBase using ESMFold; ProteinMPNN
log-likelihood scores and redesign recovery were computed against the
ESMFold backbone. Sequence and structural novelty were assessed via
mmseqs2 sequence identity and Foldseek TM-score against AFDB50.

No complex re-fold against the RBX1 target is included in the
ProteinBase release; Boltz-2, Chai-1, and Protenix-v2 complex
predictions are computed separately in this repository (figS4).

\hypertarget{statistics}{%
\subsubsection{Statistics}\label{statistics}}

We report headline counts and rates with raw n/N. Where binary
comparisons were tested (e.g.~method-family hit rates), we used Fisher's
exact (two-sided) and report the p-value alongside the odds ratio. With
9 binders across 322 designs, almost every per-method or per-modality
split is underpowered.

\hypertarget{code-and-data}{%
\subsubsection{Code and data}\label{code-and-data}}

All sequences, ESMFold predictions, and SPR sensorgrams are at
ProteinBase under ODC-ODbL. The canonical 322-row design table, complex
re-folds (Boltz-2, Chai-1, Protenix-v2), and analysis code are released
alongside this paper under MIT (code) and CC-BY-4.0 (data and figures).
